% ======================================================
% ANNEXES
% ======================================================
\newpage
\phantomsection
\addcontentsline{toc}{section}{Annexes}

\begin{center}
{\Huge\bfseries Annexes}
\end{center}

\vspace{1cm}

%================================================================
\phantomsection
\section*{Annexe A \quad Glossaire budg\'etaire}
\addcontentsline{toc}{section}{Annexe A \quad Glossaire budg\'etaire}
%================================================================

Cette annexe regroupe les principales notions budg\'etaires utilis\'ees
tout au long du m\'emoire. Les d\'efinitions retenues sont conformes au
cadre de la loi organique relative \`a la loi de finances (LOF~130-13)
et aux usages de la Direction du Budget.

\begin{description}
\item[Cr\'edits ouverts.] Montants autoris\'es par la loi de finances
pour une ligne budg\'etaire donn\'ee au titre d'un exercice. Ils
constituent le plafond th\'eorique de la d\'epense.

\item[Cr\'edits ouverts vis\'es.] Cr\'edits ouverts apr\`es prise en
compte des virements (en plus et en moins) et des r\'egulations
intervenus en cours d'exercice. C'est l'enveloppe r\'eellement
disponible pour engager la d\'epense.

\item[Engagement.] Acte par lequel l'administration cr\'ee ou constate
\`a son encontre une obligation de laquelle r\'esultera une d\'epense.
Il pr\'ec\`ede la liquidation et le paiement effectif.

\item[Total engag\'e vis\'e.] Cumul des engagements valid\'es au titre
de la ligne budg\'etaire \`a la date d'extraction du fichier source.

\item[Ordonnancement.] Acte par lequel l'ordonnateur donne l'ordre au
comptable de payer la d\'epense ayant fait l'objet d'un engagement
liquid\'e.

\item[Taux d'engagement.] Rapport, exprim\'e en pourcentage, entre le
total engag\'e et les cr\'edits ouverts vis\'es. Il mesure le degr\'e
de consommation des cr\'edits sur l'exercice. C'est l'indicateur
central exploit\'e par le dispositif propos\'e.

\item[Virement de cr\'edits.] Mouvement de cr\'edits d'une ligne
budg\'etaire vers une autre, op\'er\'e en cours d'exercice dans les
limites pr\'evues par la r\'eglementation. Les virements en plus
(\og Virements~+~\fg{}) augmentent les cr\'edits ouverts vis\'es de la
ligne b\'en\'eficiaire~; les virements en moins
(\og Virements~$-$~\fg{}) les diminuent.

\item[R\'egulation budg\'etaire.] Mesure prise en cours d'exercice par
le Ministre charg\'e des finances pour ajuster le rythme de la d\'epense
aux contraintes de tr\'esorerie ou \`a l'\'evolution des recettes.

\item[Chapitre.] Premier niveau de la nomenclature budg\'etaire.
Identifie une grande cat\'egorie de d\'epenses (personnel,
fonctionnement, investissement) au sein d'un d\'epartement minist\'eriel.

\item[Programme.] Second niveau de la nomenclature, regroupant un
ensemble coh\'erent d'actions concourant \`a une politique publique
identifi\'ee, conform\'ement \`a la d\'emarche introduite par la
LOF~130-13.

\item[R\'egion.] Niveau de la nomenclature traduisant la
d\'econcentration territoriale des cr\'edits, dans le cadre de la
r\'egionalisation avanc\'ee.

\item[Projet ou action.] Niveau interm\'ediaire de la nomenclature,
correspondant \`a une op\'eration ou \`a une activit\'e d\'eclin\'ee
au sein d'un programme.

\item[Ligne budg\'etaire.] Niveau le plus fin de la nomenclature,
identifi\'e par la combinaison Chapitre--Programme--R\'egion--Projet
--Ligne. Constitue l'unit\'e d'analyse du dispositif (107 lignes
retenues sur le p\'erim\`etre du Minist\`ere de la Justice).

\item[Cl\^oture de l'exercice.] Op\'eration comptable de fin d'ann\'ee
qui arr\^ete d\'efinitivement les engagements et les paiements de
l'exercice consid\'er\'e. Les donn\'ees relatives aux exercices
cl\^otur\'es sont consid\'er\'ees comme stables et exploitables \`a
des fins d'analyse historique.
\end{description}

\newpage

%================================================================
\phantomsection
\section*{Annexe B \quad Format des fichiers sources TGR}
\addcontentsline{toc}{section}{Annexe B \quad Format des fichiers sources TGR}
%================================================================

Le dispositif s'appuie sur les fichiers \texttt{SituationChap-AAAA.xls}
produits annuellement par la Tr\'esorerie G\'en\'erale du Royaume \`a
partir du syst\`eme d'information budg\'etaire. Chaque fichier d\'ecrit
la situation des chapitres pour l'exercice consid\'er\'e (\texttt{AAAA}
d\'esignant l'ann\'ee).

Le tableau~\ref{tab:annexe-format-tgr} pr\'ecise les colonnes
effectivement exploit\'ees par le module d'importation, leur position
dans le classeur (lettre de colonne au sens Excel), leur intitul\'e et
leur usage dans les traitements en aval. Les donn\'ees commencent
\`a la ligne~6 du classeur (\texttt{DATA\_START\_ROW = 6})~; les lignes
ant\'erieures contiennent les en-t\^etes et les libell\'es g\'en\'eraux.

\begin{table}[H]
\centering
\caption{Colonnes exploit\'ees du fichier \texttt{SituationChap-AAAA.xls}.}
\label{tab:annexe-format-tgr}
\renewcommand{\arraystretch}{1.25}
\small
\begin{tabular}{@{}cllp{6.4cm}@{}}
\hline
\textbf{Col.} & \textbf{Intitul\'e source} & \textbf{Type} & \textbf{Usage dans le dispositif} \\
\hline
B  & Chapitre                        & Code num\'erique & Identifiant du niveau~1 (cl\'e ligne) \\
C  & Programme                       & Code num\'erique & Identifiant du niveau~2 (cl\'e ligne) \\
D  & R\'egion                        & Code num\'erique & Identifiant du niveau~3 (cl\'e ligne) \\
E  & Projet/Action                   & Code num\'erique & Identifiant du niveau~4 (cl\'e ligne) \\
G  & Ligne budg\'etaire              & Code num\'erique & Identifiant du niveau~5 (cl\'e finale) \\
H  & Intitul\'e                      & Texte           & Libell\'e associ\'e au niveau courant \\
K  & Cr\'edits ouverts vis\'es       & Mon\'etaire (MAD) & Enveloppe disponible (d\'enominateur du taux d'engagement) \\
M  & Virements en plus               & Mon\'etaire (MAD) & Mouvement positif sur les cr\'edits \\
O  & Virements en moins              & Mon\'etaire (MAD) & Mouvement n\'egatif sur les cr\'edits \\
AH & Total engag\'e vis\'e           & Mon\'etaire (MAD) & Consommation effective (num\'erateur du taux d'engagement) \\
\hline
\end{tabular}
\end{table}

\paragraph{Hypoth\`eses de structure.}
Le module d'importation suppose~: (i)~la stabilit\'e des positions de
colonnes \`a travers les exercices~; (ii)~le d\'ebut effectif des
donn\'ees \`a la ligne~6~; (iii)~une organisation hi\'erarchique des
codes (un code de niveau~$k$ devient nul lorsqu'un code de niveau
$k' < k$ est rencontr\'e). Ces hypoth\`eses sont v\'erifi\'ees sur
l'ensemble des exercices~2020--2025 utilis\'es. Une \'evolution du
format des fichiers sources n\'ecessiterait une adaptation du module
d'importation, comme rappel\'e en~\S\,5.6.1.

\paragraph{Lignes exclues du p\'erim\`etre.}
Deux chapitres sont syst\'ematiquement exclus du p\'erim\`etre
analys\'e car ils correspondent \`a des cr\'edits g\'er\'es selon des
r\`egles sp\'ecifiques (r\'egulations centrales, charges communes)
qui ne se pr\^etent pas \`a la mod\'elisation propos\'ee~:
\texttt{3200106001} et \texttt{3200106002}. Le p\'erim\`etre final
retenu apr\`es exclusion et apr\`es application du crit\`ere de
continuit\'e (pr\'esence sur l'ensemble des exercices) s'\'el\`eve \`a
\textbf{107~lignes budg\'etaires}.
